%% ════════════════════════════════════════════════════════════════════════════
%%  macros.tex — Definiciones reutilizables para casos de prueba
%%  Incluir en el preámbulo de main.tex con %% ════════════════════════════════════════════════════════════════════════════
%%  macros.tex — Definiciones reutilizables para casos de prueba
%%  Incluir en el preámbulo de main.tex con %% ════════════════════════════════════════════════════════════════════════════
%%  macros.tex — Definiciones reutilizables para casos de prueba
%%  Incluir en el preámbulo de main.tex con %% ════════════════════════════════════════════════════════════════════════════
%%  macros.tex — Definiciones reutilizables para casos de prueba
%%  Incluir en el preámbulo de main.tex con \input{sections/macros}
%% ════════════════════════════════════════════════════════════════════════════

%% ── Símbolo del colón costarricense ─────────────────────────────────────────
\newcommand{\CRC}{\textit{c}/}

%% ── Flecha modular ───────────────────────────────────────────────────────────
\newcommand{\arr}{\ensuremath{\rightarrow}}

%% ── Colores ──────────────────────────────────────────────────────────────────
\definecolor{headerblue}{RGB}{31,73,125}
\definecolor{grisheader}{RGB}{195,195,195}
\definecolor{rowgray}{RGB}{235,235,235}
\definecolor{verde}{RGB}{0,130,0}
\definecolor{rojo}{RGB}{185,0,0}
\definecolor{naranja}{RGB}{195,95,0}

%% ── Estados de caso de prueba ────────────────────────────────────────────────
\newcommand{\estadoAprobado}{\textcolor{verde}{\textbf{Aprobado}}}
\newcommand{\estadoFallido}{\textcolor{rojo}{\textbf{Fallido}}}
\newcommand{\estadoParcial}{\textcolor{naranja}{\textbf{Parcial}}}
\newcommand{\estadoNoImpl}{\textcolor{gray}{\textbf{No implementado}}}
\newcommand{\estadoObs}{\textcolor{blue}{\textbf{Observación}}}

%% ── Severidad de hallazgos ───────────────────────────────────────────────────
\newcommand{\sevAlta}{\textcolor{rojo}{\textbf{Alta}}}
\newcommand{\sevMedia}{\textcolor{naranja}{\textbf{Media}}}
\newcommand{\sevBaja}{\textcolor{olive}{\textbf{Baja}}}

%% ── Configuración de hyperlinks ──────────────────────────────────────────────
\hypersetup{colorlinks=true, linkcolor=headerblue, urlcolor=headerblue}

%% ── Encabezado y pie de página ───────────────────────────────────────────────
\pagestyle{fancy}
\fancyhf{}
\fancyhead[L]{\small Universidad Nacional --- Sede Región Brunca}
\fancyhead[R]{\small Ingeniería en Sistemas III}
\fancyfoot[C]{\thepage}
\renewcommand{\headrulewidth}{0.4pt}

%% ── Formato de secciones ─────────────────────────────────────────────────────
\titleformat{\section}{\large\bfseries\color{headerblue}}{}{0em}{}[\titlerule]
\titleformat{\subsection}{\normalsize\bfseries\color{headerblue}}{}{0em}{}

%% ── Espaciado global de tablas y listas ──────────────────────────────────────
\setlength{\tabcolsep}{8pt}
\renewcommand{\arraystretch}{1.65}
\setlength{\emergencystretch}{3em}
\setlist{topsep=3pt, itemsep=2pt, parsep=0pt, partopsep=0pt, leftmargin=1.8em}

%% ── Longitudes calculadas para las tablas de casos de prueba ─────────────────
%%
%%  Tabla 3 columnas: |p{CPcolA}|p{CPcolB}|p{CPcolC}|
%%    multicolumn(3) → CPmcw   = textwidth − 2×tabcolsep − 2×arrayrulewidth
%%    multicolumn(2) → CPtipoW = CPcolA + CPcolB + 2×tabcolsep + arrayrulewidth
%%    Pasos/Result   → CPhalfW ≈ CPmcw/2 − tabcolsep − arrayrulewidth/2
%%  Tabla hallazgo 2 columnas: |p{CPcolA}|p{HALcolB}|
%%
\newlength{\CPcolA}\setlength{\CPcolA}{2.4cm}
\newlength{\CPcolB}\setlength{\CPcolB}{3.3cm}
\newlength{\CPcolC}
\newlength{\CPmcw}
\newlength{\CPtipoW}
\newlength{\CPhalfW}
\newlength{\HALcolB}

\AtBeginDocument{%
  \setlength{\CPcolC}{\textwidth}%
    \addtolength{\CPcolC}{-\CPcolA}%
    \addtolength{\CPcolC}{-\CPcolB}%
    \addtolength{\CPcolC}{-6\tabcolsep}%
    \addtolength{\CPcolC}{-4\arrayrulewidth}%
  \setlength{\CPmcw}{\textwidth}%
    \addtolength{\CPmcw}{-2\tabcolsep}%
    \addtolength{\CPmcw}{-2\arrayrulewidth}%
    \addtolength{\CPmcw}{-2pt}%
  \setlength{\CPtipoW}{\CPcolA}%
    \addtolength{\CPtipoW}{\CPcolB}%
    \addtolength{\CPtipoW}{2\tabcolsep}%
    \addtolength{\CPtipoW}{\arrayrulewidth}%
  \setlength{\HALcolB}{\textwidth}%
    \addtolength{\HALcolB}{-\CPcolA}%
    \addtolength{\HALcolB}{-4\tabcolsep}%
    \addtolength{\HALcolB}{-3\arrayrulewidth}%
    \addtolength{\HALcolB}{-2pt}%
  \setlength{\CPhalfW}{0.5\CPmcw}%
    \addtolength{\CPhalfW}{-\tabcolsep}%
    \addtolength{\CPhalfW}{-0.5\arrayrulewidth}%
    \addtolength{\CPhalfW}{-3pt}%
}

%% ── Macro \casoprueba ────────────────────────────────────────────────────────
%%
%%  Definir ANTES de cada llamada:
%%    \def\cpPasos{...}       ítems \item del enumerate de pasos
%%    \def\cpResultadoE{...}  ítems \item del itemize de resultado esperado
%%    \def\cpResultadoO{...}  texto libre del resultado obtenido
%%
%%  Sintaxis:
%%    \casoprueba{ID}{Fecha}{Módulo}{Nombre}{Tipo}{Estado}
%%               {ítems precondición}{ítems datos de prueba}{Ejecutado por}
%%
\def\cpPasos{}
\def\cpResultadoE{}
\def\cpResultadoO{}

\newcommand{\casoprueba}[9]{%
  \noindent
  \begin{tabular}{|p{\CPcolA}|p{\CPcolB}|p{\CPcolC}|}
  \hline
  \rowcolor{grisheader}%
  \multicolumn{3}{|p{\CPmcw}|}{\centering\bfseries\large Caso de prueba}\\[2pt]
  \hline
  \textbf{ID:}~#1 & \textbf{Fecha:}~#2 & \textbf{Módulo:}~#3 \\
  \hline
  \multicolumn{3}{|p{\CPmcw}|}{\textbf{Nombre:}~#4}\\
  \hline
  \multicolumn{2}{|p{\CPtipoW}|}{\textbf{Tipo de Prueba:}~#5} &
    \textbf{Estado:}~#6\\
  \hline
  \multicolumn{3}{|p{\CPmcw}|}{%
    \vspace{3pt}\textbf{Precondición:}\vspace{-4pt}%
    \begin{itemize}#7\end{itemize}\vspace{2pt}%
  }\\
  \hline
  \multicolumn{3}{|p{\CPmcw}|}{%
    \vspace{3pt}\textbf{Datos de Prueba:}\vspace{-4pt}%
    \begin{itemize}#8\end{itemize}\vspace{2pt}%
  }\\
  \hline
  \multicolumn{3}{|p{\CPmcw}|}{%
    \renewcommand{\arraystretch}{1.0}%
    \begin{tabular}[t]{@{}p{\CPhalfW}|p{\CPhalfW}@{}}
      \vspace{3pt}\textbf{Pasos:}\vspace{-4pt}%
      \begin{enumerate}\cpPasos\end{enumerate}\vspace{2pt}%
      &
      \vspace{3pt}\textbf{Resultado Esperado:}\vspace{-4pt}%
      \begin{itemize}\cpResultadoE\end{itemize}\vspace{2pt}%
      \\
    \end{tabular}%
  }\\
  \hline
  \multicolumn{3}{|p{\CPmcw}|}{%
    \vspace{3pt}\textbf{Resultado Obtenido:}\par\smallskip
    \cpResultadoO\vspace{5pt}%
  }\\
  \hline
  \multicolumn{3}{|p{\CPmcw}|}{%
    \vspace{2pt}\textbf{Prueba ejecutada por:}~#9\vspace{2pt}%
  }\\
  \hline
  \end{tabular}%
  \par\vspace{1cm}%
}

%% ── Macro \hallazgo ──────────────────────────────────────────────────────────
%%
%%  Sintaxis:
%%    \hallazgo{ID}{Título}{Severidad}{Módulo}{Tipo}{Caso relacionado}
%%             {Descripción}{Impacto}{Recomendación}
%%
\newcommand{\hallazgo}[9]{%
  \noindent
  \begin{tabular}{|p{\CPcolA}|p{\HALcolB}|}
  \hline
  \rowcolor{grisheader}%
  \multicolumn{2}{|p{\CPmcw}|}{\centering\bfseries\large #1 --- #2}\\[2pt]
  \hline
  \textbf{ID}               & #1 \\
  \hline
  \textbf{Severidad}        & #3 \\
  \hline
  \textbf{Módulo}           & #4 \\
  \hline
  \textbf{Tipo}             & #5 \\
  \hline
  \textbf{Caso relacionado} & #6 \\
  \hline
  \textbf{Estado}           & Confirmado \\
  \hline
  \multicolumn{2}{|p{\CPmcw}|}{%
    \vspace{3pt}\textbf{Descripción:}\par\smallskip #7\vspace{5pt}%
  }\\
  \hline
  \multicolumn{2}{|p{\CPmcw}|}{%
    \vspace{3pt}\textbf{Impacto:}\par\smallskip #8\vspace{5pt}%
  }\\
  \hline
  \multicolumn{2}{|p{\CPmcw}|}{%
    \vspace{3pt}\textbf{Recomendación:}\par\smallskip #9\vspace{5pt}%
  }\\
  \hline
  \end{tabular}%
  \par\vspace{1cm}%
}

%% ════════════════════════════════════════════════════════════════════════════

%% ── Símbolo del colón costarricense ─────────────────────────────────────────
\newcommand{\CRC}{\textit{c}/}

%% ── Flecha modular ───────────────────────────────────────────────────────────
\newcommand{\arr}{\ensuremath{\rightarrow}}

%% ── Colores ──────────────────────────────────────────────────────────────────
\definecolor{headerblue}{RGB}{31,73,125}
\definecolor{grisheader}{RGB}{195,195,195}
\definecolor{rowgray}{RGB}{235,235,235}
\definecolor{verde}{RGB}{0,130,0}
\definecolor{rojo}{RGB}{185,0,0}
\definecolor{naranja}{RGB}{195,95,0}

%% ── Estados de caso de prueba ────────────────────────────────────────────────
\newcommand{\estadoAprobado}{\textcolor{verde}{\textbf{Aprobado}}}
\newcommand{\estadoFallido}{\textcolor{rojo}{\textbf{Fallido}}}
\newcommand{\estadoParcial}{\textcolor{naranja}{\textbf{Parcial}}}
\newcommand{\estadoNoImpl}{\textcolor{gray}{\textbf{No implementado}}}
\newcommand{\estadoObs}{\textcolor{blue}{\textbf{Observación}}}

%% ── Severidad de hallazgos ───────────────────────────────────────────────────
\newcommand{\sevAlta}{\textcolor{rojo}{\textbf{Alta}}}
\newcommand{\sevMedia}{\textcolor{naranja}{\textbf{Media}}}
\newcommand{\sevBaja}{\textcolor{olive}{\textbf{Baja}}}

%% ── Configuración de hyperlinks ──────────────────────────────────────────────
\hypersetup{colorlinks=true, linkcolor=headerblue, urlcolor=headerblue}

%% ── Encabezado y pie de página ───────────────────────────────────────────────
\pagestyle{fancy}
\fancyhf{}
\fancyhead[L]{\small Universidad Nacional --- Sede Región Brunca}
\fancyhead[R]{\small Ingeniería en Sistemas III}
\fancyfoot[C]{\thepage}
\renewcommand{\headrulewidth}{0.4pt}

%% ── Formato de secciones ─────────────────────────────────────────────────────
\titleformat{\section}{\large\bfseries\color{headerblue}}{}{0em}{}[\titlerule]
\titleformat{\subsection}{\normalsize\bfseries\color{headerblue}}{}{0em}{}

%% ── Espaciado global de tablas y listas ──────────────────────────────────────
\setlength{\tabcolsep}{8pt}
\renewcommand{\arraystretch}{1.65}
\setlength{\emergencystretch}{3em}
\setlist{topsep=3pt, itemsep=2pt, parsep=0pt, partopsep=0pt, leftmargin=1.8em}

%% ── Longitudes calculadas para las tablas de casos de prueba ─────────────────
%%
%%  Tabla 3 columnas: |p{CPcolA}|p{CPcolB}|p{CPcolC}|
%%    multicolumn(3) → CPmcw   = textwidth − 2×tabcolsep − 2×arrayrulewidth
%%    multicolumn(2) → CPtipoW = CPcolA + CPcolB + 2×tabcolsep + arrayrulewidth
%%    Pasos/Result   → CPhalfW ≈ CPmcw/2 − tabcolsep − arrayrulewidth/2
%%  Tabla hallazgo 2 columnas: |p{CPcolA}|p{HALcolB}|
%%
\newlength{\CPcolA}\setlength{\CPcolA}{2.4cm}
\newlength{\CPcolB}\setlength{\CPcolB}{3.3cm}
\newlength{\CPcolC}
\newlength{\CPmcw}
\newlength{\CPtipoW}
\newlength{\CPhalfW}
\newlength{\HALcolB}

\AtBeginDocument{%
  \setlength{\CPcolC}{\textwidth}%
    \addtolength{\CPcolC}{-\CPcolA}%
    \addtolength{\CPcolC}{-\CPcolB}%
    \addtolength{\CPcolC}{-6\tabcolsep}%
    \addtolength{\CPcolC}{-4\arrayrulewidth}%
  \setlength{\CPmcw}{\textwidth}%
    \addtolength{\CPmcw}{-2\tabcolsep}%
    \addtolength{\CPmcw}{-2\arrayrulewidth}%
    \addtolength{\CPmcw}{-2pt}%
  \setlength{\CPtipoW}{\CPcolA}%
    \addtolength{\CPtipoW}{\CPcolB}%
    \addtolength{\CPtipoW}{2\tabcolsep}%
    \addtolength{\CPtipoW}{\arrayrulewidth}%
  \setlength{\HALcolB}{\textwidth}%
    \addtolength{\HALcolB}{-\CPcolA}%
    \addtolength{\HALcolB}{-4\tabcolsep}%
    \addtolength{\HALcolB}{-3\arrayrulewidth}%
    \addtolength{\HALcolB}{-2pt}%
  \setlength{\CPhalfW}{0.5\CPmcw}%
    \addtolength{\CPhalfW}{-\tabcolsep}%
    \addtolength{\CPhalfW}{-0.5\arrayrulewidth}%
    \addtolength{\CPhalfW}{-3pt}%
}

%% ── Macro \casoprueba ────────────────────────────────────────────────────────
%%
%%  Definir ANTES de cada llamada:
%%    \def\cpPasos{...}       ítems \item del enumerate de pasos
%%    \def\cpResultadoE{...}  ítems \item del itemize de resultado esperado
%%    \def\cpResultadoO{...}  texto libre del resultado obtenido
%%
%%  Sintaxis:
%%    \casoprueba{ID}{Fecha}{Módulo}{Nombre}{Tipo}{Estado}
%%               {ítems precondición}{ítems datos de prueba}{Ejecutado por}
%%
\def\cpPasos{}
\def\cpResultadoE{}
\def\cpResultadoO{}

\newcommand{\casoprueba}[9]{%
  \noindent
  \begin{tabular}{|p{\CPcolA}|p{\CPcolB}|p{\CPcolC}|}
  \hline
  \rowcolor{grisheader}%
  \multicolumn{3}{|p{\CPmcw}|}{\centering\bfseries\large Caso de prueba}\\[2pt]
  \hline
  \textbf{ID:}~#1 & \textbf{Fecha:}~#2 & \textbf{Módulo:}~#3 \\
  \hline
  \multicolumn{3}{|p{\CPmcw}|}{\textbf{Nombre:}~#4}\\
  \hline
  \multicolumn{2}{|p{\CPtipoW}|}{\textbf{Tipo de Prueba:}~#5} &
    \textbf{Estado:}~#6\\
  \hline
  \multicolumn{3}{|p{\CPmcw}|}{%
    \vspace{3pt}\textbf{Precondición:}\vspace{-4pt}%
    \begin{itemize}#7\end{itemize}\vspace{2pt}%
  }\\
  \hline
  \multicolumn{3}{|p{\CPmcw}|}{%
    \vspace{3pt}\textbf{Datos de Prueba:}\vspace{-4pt}%
    \begin{itemize}#8\end{itemize}\vspace{2pt}%
  }\\
  \hline
  \multicolumn{3}{|p{\CPmcw}|}{%
    \renewcommand{\arraystretch}{1.0}%
    \begin{tabular}[t]{@{}p{\CPhalfW}|p{\CPhalfW}@{}}
      \vspace{3pt}\textbf{Pasos:}\vspace{-4pt}%
      \begin{enumerate}\cpPasos\end{enumerate}\vspace{2pt}%
      &
      \vspace{3pt}\textbf{Resultado Esperado:}\vspace{-4pt}%
      \begin{itemize}\cpResultadoE\end{itemize}\vspace{2pt}%
      \\
    \end{tabular}%
  }\\
  \hline
  \multicolumn{3}{|p{\CPmcw}|}{%
    \vspace{3pt}\textbf{Resultado Obtenido:}\par\smallskip
    \cpResultadoO\vspace{5pt}%
  }\\
  \hline
  \multicolumn{3}{|p{\CPmcw}|}{%
    \vspace{2pt}\textbf{Prueba ejecutada por:}~#9\vspace{2pt}%
  }\\
  \hline
  \end{tabular}%
  \par\vspace{1cm}%
}

%% ── Macro \hallazgo ──────────────────────────────────────────────────────────
%%
%%  Sintaxis:
%%    \hallazgo{ID}{Título}{Severidad}{Módulo}{Tipo}{Caso relacionado}
%%             {Descripción}{Impacto}{Recomendación}
%%
\newcommand{\hallazgo}[9]{%
  \noindent
  \begin{tabular}{|p{\CPcolA}|p{\HALcolB}|}
  \hline
  \rowcolor{grisheader}%
  \multicolumn{2}{|p{\CPmcw}|}{\centering\bfseries\large #1 --- #2}\\[2pt]
  \hline
  \textbf{ID}               & #1 \\
  \hline
  \textbf{Severidad}        & #3 \\
  \hline
  \textbf{Módulo}           & #4 \\
  \hline
  \textbf{Tipo}             & #5 \\
  \hline
  \textbf{Caso relacionado} & #6 \\
  \hline
  \textbf{Estado}           & Confirmado \\
  \hline
  \multicolumn{2}{|p{\CPmcw}|}{%
    \vspace{3pt}\textbf{Descripción:}\par\smallskip #7\vspace{5pt}%
  }\\
  \hline
  \multicolumn{2}{|p{\CPmcw}|}{%
    \vspace{3pt}\textbf{Impacto:}\par\smallskip #8\vspace{5pt}%
  }\\
  \hline
  \multicolumn{2}{|p{\CPmcw}|}{%
    \vspace{3pt}\textbf{Recomendación:}\par\smallskip #9\vspace{5pt}%
  }\\
  \hline
  \end{tabular}%
  \par\vspace{1cm}%
}

%% ════════════════════════════════════════════════════════════════════════════

%% ── Símbolo del colón costarricense ─────────────────────────────────────────
\newcommand{\CRC}{\textit{c}/}

%% ── Flecha modular ───────────────────────────────────────────────────────────
\newcommand{\arr}{\ensuremath{\rightarrow}}

%% ── Colores ──────────────────────────────────────────────────────────────────
\definecolor{headerblue}{RGB}{31,73,125}
\definecolor{grisheader}{RGB}{195,195,195}
\definecolor{rowgray}{RGB}{235,235,235}
\definecolor{verde}{RGB}{0,130,0}
\definecolor{rojo}{RGB}{185,0,0}
\definecolor{naranja}{RGB}{195,95,0}

%% ── Estados de caso de prueba ────────────────────────────────────────────────
\newcommand{\estadoAprobado}{\textcolor{verde}{\textbf{Aprobado}}}
\newcommand{\estadoFallido}{\textcolor{rojo}{\textbf{Fallido}}}
\newcommand{\estadoParcial}{\textcolor{naranja}{\textbf{Parcial}}}
\newcommand{\estadoNoImpl}{\textcolor{gray}{\textbf{No implementado}}}
\newcommand{\estadoObs}{\textcolor{blue}{\textbf{Observación}}}

%% ── Severidad de hallazgos ───────────────────────────────────────────────────
\newcommand{\sevAlta}{\textcolor{rojo}{\textbf{Alta}}}
\newcommand{\sevMedia}{\textcolor{naranja}{\textbf{Media}}}
\newcommand{\sevBaja}{\textcolor{olive}{\textbf{Baja}}}

%% ── Configuración de hyperlinks ──────────────────────────────────────────────
\hypersetup{colorlinks=true, linkcolor=headerblue, urlcolor=headerblue}

%% ── Encabezado y pie de página ───────────────────────────────────────────────
\pagestyle{fancy}
\fancyhf{}
\fancyhead[L]{\small Universidad Nacional --- Sede Región Brunca}
\fancyhead[R]{\small Ingeniería en Sistemas III}
\fancyfoot[C]{\thepage}
\renewcommand{\headrulewidth}{0.4pt}

%% ── Formato de secciones ─────────────────────────────────────────────────────
\titleformat{\section}{\large\bfseries\color{headerblue}}{}{0em}{}[\titlerule]
\titleformat{\subsection}{\normalsize\bfseries\color{headerblue}}{}{0em}{}

%% ── Espaciado global de tablas y listas ──────────────────────────────────────
\setlength{\tabcolsep}{8pt}
\renewcommand{\arraystretch}{1.65}
\setlength{\emergencystretch}{3em}
\setlist{topsep=3pt, itemsep=2pt, parsep=0pt, partopsep=0pt, leftmargin=1.8em}

%% ── Longitudes calculadas para las tablas de casos de prueba ─────────────────
%%
%%  Tabla 3 columnas: |p{CPcolA}|p{CPcolB}|p{CPcolC}|
%%    multicolumn(3) → CPmcw   = textwidth − 2×tabcolsep − 2×arrayrulewidth
%%    multicolumn(2) → CPtipoW = CPcolA + CPcolB + 2×tabcolsep + arrayrulewidth
%%    Pasos/Result   → CPhalfW ≈ CPmcw/2 − tabcolsep − arrayrulewidth/2
%%  Tabla hallazgo 2 columnas: |p{CPcolA}|p{HALcolB}|
%%
\newlength{\CPcolA}\setlength{\CPcolA}{2.4cm}
\newlength{\CPcolB}\setlength{\CPcolB}{3.3cm}
\newlength{\CPcolC}
\newlength{\CPmcw}
\newlength{\CPtipoW}
\newlength{\CPhalfW}
\newlength{\HALcolB}

\AtBeginDocument{%
  \setlength{\CPcolC}{\textwidth}%
    \addtolength{\CPcolC}{-\CPcolA}%
    \addtolength{\CPcolC}{-\CPcolB}%
    \addtolength{\CPcolC}{-6\tabcolsep}%
    \addtolength{\CPcolC}{-4\arrayrulewidth}%
  \setlength{\CPmcw}{\textwidth}%
    \addtolength{\CPmcw}{-2\tabcolsep}%
    \addtolength{\CPmcw}{-2\arrayrulewidth}%
    \addtolength{\CPmcw}{-2pt}%
  \setlength{\CPtipoW}{\CPcolA}%
    \addtolength{\CPtipoW}{\CPcolB}%
    \addtolength{\CPtipoW}{2\tabcolsep}%
    \addtolength{\CPtipoW}{\arrayrulewidth}%
  \setlength{\HALcolB}{\textwidth}%
    \addtolength{\HALcolB}{-\CPcolA}%
    \addtolength{\HALcolB}{-4\tabcolsep}%
    \addtolength{\HALcolB}{-3\arrayrulewidth}%
    \addtolength{\HALcolB}{-2pt}%
  \setlength{\CPhalfW}{0.5\CPmcw}%
    \addtolength{\CPhalfW}{-\tabcolsep}%
    \addtolength{\CPhalfW}{-0.5\arrayrulewidth}%
    \addtolength{\CPhalfW}{-3pt}%
}

%% ── Macro \casoprueba ────────────────────────────────────────────────────────
%%
%%  Definir ANTES de cada llamada:
%%    \def\cpPasos{...}       ítems \item del enumerate de pasos
%%    \def\cpResultadoE{...}  ítems \item del itemize de resultado esperado
%%    \def\cpResultadoO{...}  texto libre del resultado obtenido
%%
%%  Sintaxis:
%%    \casoprueba{ID}{Fecha}{Módulo}{Nombre}{Tipo}{Estado}
%%               {ítems precondición}{ítems datos de prueba}{Ejecutado por}
%%
\def\cpPasos{}
\def\cpResultadoE{}
\def\cpResultadoO{}

\newcommand{\casoprueba}[9]{%
  \noindent
  \begin{tabular}{|p{\CPcolA}|p{\CPcolB}|p{\CPcolC}|}
  \hline
  \rowcolor{grisheader}%
  \multicolumn{3}{|p{\CPmcw}|}{\centering\bfseries\large Caso de prueba}\\[2pt]
  \hline
  \textbf{ID:}~#1 & \textbf{Fecha:}~#2 & \textbf{Módulo:}~#3 \\
  \hline
  \multicolumn{3}{|p{\CPmcw}|}{\textbf{Nombre:}~#4}\\
  \hline
  \multicolumn{2}{|p{\CPtipoW}|}{\textbf{Tipo de Prueba:}~#5} &
    \textbf{Estado:}~#6\\
  \hline
  \multicolumn{3}{|p{\CPmcw}|}{%
    \vspace{3pt}\textbf{Precondición:}\vspace{-4pt}%
    \begin{itemize}#7\end{itemize}\vspace{2pt}%
  }\\
  \hline
  \multicolumn{3}{|p{\CPmcw}|}{%
    \vspace{3pt}\textbf{Datos de Prueba:}\vspace{-4pt}%
    \begin{itemize}#8\end{itemize}\vspace{2pt}%
  }\\
  \hline
  \multicolumn{3}{|p{\CPmcw}|}{%
    \renewcommand{\arraystretch}{1.0}%
    \begin{tabular}[t]{@{}p{\CPhalfW}|p{\CPhalfW}@{}}
      \vspace{3pt}\textbf{Pasos:}\vspace{-4pt}%
      \begin{enumerate}\cpPasos\end{enumerate}\vspace{2pt}%
      &
      \vspace{3pt}\textbf{Resultado Esperado:}\vspace{-4pt}%
      \begin{itemize}\cpResultadoE\end{itemize}\vspace{2pt}%
      \\
    \end{tabular}%
  }\\
  \hline
  \multicolumn{3}{|p{\CPmcw}|}{%
    \vspace{3pt}\textbf{Resultado Obtenido:}\par\smallskip
    \cpResultadoO\vspace{5pt}%
  }\\
  \hline
  \multicolumn{3}{|p{\CPmcw}|}{%
    \vspace{2pt}\textbf{Prueba ejecutada por:}~#9\vspace{2pt}%
  }\\
  \hline
  \end{tabular}%
  \par\vspace{1cm}%
}

%% ── Macro \hallazgo ──────────────────────────────────────────────────────────
%%
%%  Sintaxis:
%%    \hallazgo{ID}{Título}{Severidad}{Módulo}{Tipo}{Caso relacionado}
%%             {Descripción}{Impacto}{Recomendación}
%%
\newcommand{\hallazgo}[9]{%
  \noindent
  \begin{tabular}{|p{\CPcolA}|p{\HALcolB}|}
  \hline
  \rowcolor{grisheader}%
  \multicolumn{2}{|p{\CPmcw}|}{\centering\bfseries\large #1 --- #2}\\[2pt]
  \hline
  \textbf{ID}               & #1 \\
  \hline
  \textbf{Severidad}        & #3 \\
  \hline
  \textbf{Módulo}           & #4 \\
  \hline
  \textbf{Tipo}             & #5 \\
  \hline
  \textbf{Caso relacionado} & #6 \\
  \hline
  \textbf{Estado}           & Confirmado \\
  \hline
  \multicolumn{2}{|p{\CPmcw}|}{%
    \vspace{3pt}\textbf{Descripción:}\par\smallskip #7\vspace{5pt}%
  }\\
  \hline
  \multicolumn{2}{|p{\CPmcw}|}{%
    \vspace{3pt}\textbf{Impacto:}\par\smallskip #8\vspace{5pt}%
  }\\
  \hline
  \multicolumn{2}{|p{\CPmcw}|}{%
    \vspace{3pt}\textbf{Recomendación:}\par\smallskip #9\vspace{5pt}%
  }\\
  \hline
  \end{tabular}%
  \par\vspace{1cm}%
}

%% ════════════════════════════════════════════════════════════════════════════

%% ── Símbolo del colón costarricense ─────────────────────────────────────────
\newcommand{\CRC}{\textit{c}/}

%% ── Flecha modular ───────────────────────────────────────────────────────────
\newcommand{\arr}{\ensuremath{\rightarrow}}

%% ── Colores ──────────────────────────────────────────────────────────────────
\definecolor{headerblue}{RGB}{31,73,125}
\definecolor{grisheader}{RGB}{195,195,195}
\definecolor{rowgray}{RGB}{235,235,235}
\definecolor{verde}{RGB}{0,130,0}
\definecolor{rojo}{RGB}{185,0,0}
\definecolor{naranja}{RGB}{195,95,0}

%% ── Estados de caso de prueba ────────────────────────────────────────────────
\newcommand{\estadoAprobado}{\textcolor{verde}{\textbf{Aprobado}}}
\newcommand{\estadoFallido}{\textcolor{rojo}{\textbf{Fallido}}}
\newcommand{\estadoParcial}{\textcolor{naranja}{\textbf{Parcial}}}
\newcommand{\estadoNoImpl}{\textcolor{gray}{\textbf{No implementado}}}
\newcommand{\estadoObs}{\textcolor{blue}{\textbf{Observación}}}

%% ── Severidad de hallazgos ───────────────────────────────────────────────────
\newcommand{\sevAlta}{\textcolor{rojo}{\textbf{Alta}}}
\newcommand{\sevMedia}{\textcolor{naranja}{\textbf{Media}}}
\newcommand{\sevBaja}{\textcolor{olive}{\textbf{Baja}}}

%% ── Configuración de hyperlinks ──────────────────────────────────────────────
\hypersetup{colorlinks=true, linkcolor=headerblue, urlcolor=headerblue}

%% ── Encabezado y pie de página ───────────────────────────────────────────────
\pagestyle{fancy}
\fancyhf{}
\fancyhead[L]{\small Universidad Nacional --- Sede Región Brunca}
\fancyhead[R]{\small Ingeniería en Sistemas III}
\fancyfoot[C]{\thepage}
\renewcommand{\headrulewidth}{0.4pt}

%% ── Formato de secciones ─────────────────────────────────────────────────────
\titleformat{\section}{\large\bfseries\color{headerblue}}{}{0em}{}[\titlerule]
\titleformat{\subsection}{\normalsize\bfseries\color{headerblue}}{}{0em}{}

%% ── Espaciado global de tablas y listas ──────────────────────────────────────
\setlength{\tabcolsep}{8pt}
\renewcommand{\arraystretch}{1.65}
\setlength{\emergencystretch}{3em}
\setlist{topsep=3pt, itemsep=2pt, parsep=0pt, partopsep=0pt, leftmargin=1.8em}

%% ── Longitudes calculadas para las tablas de casos de prueba ─────────────────
%%
%%  Tabla 3 columnas: |p{CPcolA}|p{CPcolB}|p{CPcolC}|
%%    multicolumn(3) → CPmcw   = textwidth − 2×tabcolsep − 2×arrayrulewidth
%%    multicolumn(2) → CPtipoW = CPcolA + CPcolB + 2×tabcolsep + arrayrulewidth
%%    Pasos/Result   → CPhalfW ≈ CPmcw/2 − tabcolsep − arrayrulewidth/2
%%  Tabla hallazgo 2 columnas: |p{CPcolA}|p{HALcolB}|
%%
\newlength{\CPcolA}\setlength{\CPcolA}{2.4cm}
\newlength{\CPcolB}\setlength{\CPcolB}{3.3cm}
\newlength{\CPcolC}
\newlength{\CPmcw}
\newlength{\CPtipoW}
\newlength{\CPhalfW}
\newlength{\HALcolB}

\AtBeginDocument{%
  \setlength{\CPcolC}{\textwidth}%
    \addtolength{\CPcolC}{-\CPcolA}%
    \addtolength{\CPcolC}{-\CPcolB}%
    \addtolength{\CPcolC}{-6\tabcolsep}%
    \addtolength{\CPcolC}{-4\arrayrulewidth}%
  \setlength{\CPmcw}{\textwidth}%
    \addtolength{\CPmcw}{-2\tabcolsep}%
    \addtolength{\CPmcw}{-2\arrayrulewidth}%
    \addtolength{\CPmcw}{-2pt}%
  \setlength{\CPtipoW}{\CPcolA}%
    \addtolength{\CPtipoW}{\CPcolB}%
    \addtolength{\CPtipoW}{2\tabcolsep}%
    \addtolength{\CPtipoW}{\arrayrulewidth}%
  \setlength{\HALcolB}{\textwidth}%
    \addtolength{\HALcolB}{-\CPcolA}%
    \addtolength{\HALcolB}{-4\tabcolsep}%
    \addtolength{\HALcolB}{-3\arrayrulewidth}%
    \addtolength{\HALcolB}{-2pt}%
  \setlength{\CPhalfW}{0.5\CPmcw}%
    \addtolength{\CPhalfW}{-\tabcolsep}%
    \addtolength{\CPhalfW}{-0.5\arrayrulewidth}%
    \addtolength{\CPhalfW}{-3pt}%
}

%% ── Macro \casoprueba ────────────────────────────────────────────────────────
%%
%%  Definir ANTES de cada llamada:
%%    \def\cpPasos{...}       ítems \item del enumerate de pasos
%%    \def\cpResultadoE{...}  ítems \item del itemize de resultado esperado
%%    \def\cpResultadoO{...}  texto libre del resultado obtenido
%%
%%  Sintaxis:
%%    \casoprueba{ID}{Fecha}{Módulo}{Nombre}{Tipo}{Estado}
%%               {ítems precondición}{ítems datos de prueba}{Ejecutado por}
%%
\def\cpPasos{}
\def\cpResultadoE{}
\def\cpResultadoO{}

\newcommand{\casoprueba}[9]{%
  \noindent
  \begin{tabular}{|p{\CPcolA}|p{\CPcolB}|p{\CPcolC}|}
  \hline
  \rowcolor{grisheader}%
  \multicolumn{3}{|p{\CPmcw}|}{\centering\bfseries\large Caso de prueba}\\[2pt]
  \hline
  \textbf{ID:}~#1 & \textbf{Fecha:}~#2 & \textbf{Módulo:}~#3 \\
  \hline
  \multicolumn{3}{|p{\CPmcw}|}{\textbf{Nombre:}~#4}\\
  \hline
  \multicolumn{2}{|p{\CPtipoW}|}{\textbf{Tipo de Prueba:}~#5} &
    \textbf{Estado:}~#6\\
  \hline
  \multicolumn{3}{|p{\CPmcw}|}{%
    \vspace{3pt}\textbf{Precondición:}\vspace{-4pt}%
    \begin{itemize}#7\end{itemize}\vspace{2pt}%
  }\\
  \hline
  \multicolumn{3}{|p{\CPmcw}|}{%
    \vspace{3pt}\textbf{Datos de Prueba:}\vspace{-4pt}%
    \begin{itemize}#8\end{itemize}\vspace{2pt}%
  }\\
  \hline
  \multicolumn{3}{|p{\CPmcw}|}{%
    \renewcommand{\arraystretch}{1.0}%
    \begin{tabular}[t]{@{}p{\CPhalfW}|p{\CPhalfW}@{}}
      \vspace{3pt}\textbf{Pasos:}\vspace{-4pt}%
      \begin{enumerate}\cpPasos\end{enumerate}\vspace{2pt}%
      &
      \vspace{3pt}\textbf{Resultado Esperado:}\vspace{-4pt}%
      \begin{itemize}\cpResultadoE\end{itemize}\vspace{2pt}%
      \\
    \end{tabular}%
  }\\
  \hline
  \multicolumn{3}{|p{\CPmcw}|}{%
    \vspace{3pt}\textbf{Resultado Obtenido:}\par\smallskip
    \cpResultadoO\vspace{5pt}%
  }\\
  \hline
  \multicolumn{3}{|p{\CPmcw}|}{%
    \vspace{2pt}\textbf{Prueba ejecutada por:}~#9\vspace{2pt}%
  }\\
  \hline
  \end{tabular}%
  \par\vspace{1cm}%
}

%% ── Macro \hallazgo ──────────────────────────────────────────────────────────
%%
%%  Sintaxis:
%%    \hallazgo{ID}{Título}{Severidad}{Módulo}{Tipo}{Caso relacionado}
%%             {Descripción}{Impacto}{Recomendación}
%%
\newcommand{\hallazgo}[9]{%
  \noindent
  \begin{tabular}{|p{\CPcolA}|p{\HALcolB}|}
  \hline
  \rowcolor{grisheader}%
  \multicolumn{2}{|p{\CPmcw}|}{\centering\bfseries\large #1 --- #2}\\[2pt]
  \hline
  \textbf{ID}               & #1 \\
  \hline
  \textbf{Severidad}        & #3 \\
  \hline
  \textbf{Módulo}           & #4 \\
  \hline
  \textbf{Tipo}             & #5 \\
  \hline
  \textbf{Caso relacionado} & #6 \\
  \hline
  \textbf{Estado}           & Confirmado \\
  \hline
  \multicolumn{2}{|p{\CPmcw}|}{%
    \vspace{3pt}\textbf{Descripción:}\par\smallskip #7\vspace{5pt}%
  }\\
  \hline
  \multicolumn{2}{|p{\CPmcw}|}{%
    \vspace{3pt}\textbf{Impacto:}\par\smallskip #8\vspace{5pt}%
  }\\
  \hline
  \multicolumn{2}{|p{\CPmcw}|}{%
    \vspace{3pt}\textbf{Recomendación:}\par\smallskip #9\vspace{5pt}%
  }\\
  \hline
  \end{tabular}%
  \par\vspace{1cm}%
}
